\documentclass[tikz,border=10pt]{standalone}
\usepackage{amssymb}
\usepackage{fontspec}
\usepackage{xeCJK}

% 严格遵循字体规范：英文 Times New Roman，中文黑体
\setmainfont{Times New Roman}
\setCJKmainfont{SimHei}

\usepackage{tikz}
\usetikzlibrary{positioning, calc, backgrounds, fit, arrows.meta}

% ============================================
% IEEE 顶刊风格配色定义
% ============================================
\definecolor{serialGray}{HTML}{BDC3C7}   % 串行基准色-中性灰
\definecolor{pipeBlue}{HTML}{2874A6}     % 流水线优化色-学术钢蓝
\definecolor{highlightRed}{HTML}{B03A2E} % 突出强调色-学术砖红
\definecolor{textDark}{HTML}{2C3E50}     % 深色文字
\definecolor{textGray}{HTML}{5D6D7E}     % 灰色文字
\definecolor{borderColor}{HTML}{95A5A6}  % 边框色

\begin{document}
\begin{tikzpicture}[
  % --- 统一字号规范管理 ---
  fontCaption/.style={font=\fontsize{9pt}{10.8pt}\selectfont\bfseries}, 
  fontAxis/.style={font=\fontsize{8pt}{9.6pt}\selectfont\bfseries},    
  fontNormal/.style={font=\fontsize{7.5pt}{9pt}\selectfont},           
  fontSmall/.style={font=\fontsize{6.5pt}{7.8pt}\selectfont},
  % 坐标轴缩放：X轴 1单位 = 1.3cm, Y轴 1单位(ms) = 0.015cm
  x=1.3cm, y=0.015cm  
]

% ============================================
% 整体背景
% ============================================
\fill[white] (-1.2, -60) rectangle (9.2, 500);

% ============================================
% 标题与图例区
% ============================================
% 图表主标题
\node[fontCaption, textDark] at (4.0, 480) {异构大小核流水线性能对比 (OpenAMP 三核统一板态)};

% 图例
\fill[serialGray] (2.0, 435) rectangle (2.3, 450);
\node[right, fontNormal, textDark] at (2.4, 442.5) {串行重建 (Baseline)};

\fill[pipeBlue] (4.5, 435) rectangle (4.8, 450);
\node[right, fontNormal, textDark] at (4.9, 442.5) {流水线重建 (Optimized)};

% ============================================
% 坐标系与网格线
% ============================================
% Y轴网格线与刻度
\foreach \y in {0, 100, 200, 300, 400} {
  \draw[dashed, textGray!40, line width=0.5pt] (0, \y) -- (8, \y);
  \node[left, fontSmall, textGray] at (-0.1, \y) {\y};
}

% X轴与Y轴主线
\draw[thick, borderColor] (0,0) -- (0, 420);
\draw[thick, borderColor] (0,0) -- (8, 0);

% 轴标题
\node[fontAxis, textDark, rotate=90] at (-1.0, 200) {单张重建时间 (ms/image)};
\node[fontAxis, textDark, align=center] at (2.5, -30) {MetaSchedule 优化版本};
\node[fontAxis, textDark, align=center] at (5.5, -30) {手写优化版本};

% ============================================
% 数据柱状图与标注：左侧组 (MetaSchedule)
% ============================================
% 串行柱 (347.3 ms)
\fill[serialGray] (2.0, 0) rectangle (2.45, 347.3);
\node[above, fontAxis, textDark] at (2.225, 347.3) {347.3};

% 流水线柱 (257.4 ms)
\fill[pipeBlue] (2.55, 0) rectangle (3.0, 257.4);
\node[above, fontAxis, text=pipeBlue] at (2.775, 257.4) {257.4};

% 提升比例趋势箭头
\draw[-{Stealth[length=2.5mm, width=1.8mm]}, thick, highlightRed] 
  (2.225, 370) -- (2.225, 405) -- (2.775, 405) -- (2.775, 280);
\node[above, fontAxis, text=highlightRed] at (2.5, 405) {+34.6\% 吞吐提升};


% ============================================
% 数据柱状图与标注：右侧组 (手写优化)
% ============================================
% 串行柱 (345.6 ms)
\fill[serialGray] (5.0, 0) rectangle (5.45, 345.6);
\node[above, fontAxis, textDark] at (5.225, 345.6) {345.6};

% 流水线柱 (249.4 ms)
\fill[pipeBlue] (5.55, 0) rectangle (6.0, 249.4);
\node[above, fontAxis, text=pipeBlue] at (5.775, 249.4) {249.4};

% 提升比例趋势箭头
\draw[-{Stealth[length=2.5mm, width=1.8mm]}, thick, highlightRed] 
  (5.225, 370) -- (5.225, 405) -- (5.775, 405) -- (5.775, 275);
\node[above, fontAxis, text=highlightRed] at (5.5, 405) {+38.7\% 吞吐提升};


% ============================================
% 跨组横向对比参考线 (直观印证手写优化全面占优)
% ============================================
\draw[densely dotted, thick, highlightRed] (2.775, 249.4) -- (5.55, 249.4);
\node[below, fontSmall, text=highlightRed] at (4.16, 249.4) {249.4 < 257.4 (手写优化延迟更低)};

\end{tikzpicture}
\end{document}